\section{Supporting audit and validation studies}
\label{app:audit}
This appendix details the execution and feedback validation checks, defect recall
benchmarks, and empirical audits across model tiers.

\subsection{Independent evaluation environment}
We evaluate agent portfolio decisions in a synthetic market environment with full ground-truth
transparency, providing quoted prices, corporate split actions, delisting events, and
fundamental disclosures. A submitted strategy script returns portfolio weights, alongside
a written report stating its expected Sharpe ratio. Session-$t$ portfolio weights must use
only information available by the close of session $t-1$.

An independent oracle evaluator verifies submissions by perturbing future data, checking
for post-delisting exposure, and recalculating performance under fixed transaction costs.
Oracle grades are generated after the submission receipt is saved and are never exposed
to the agent during its decision turns.

This oracle separates genuine investment skill from data leakage. In this environment, an
honest strategy trading strictly lagged features achieves an empirical ceiling Sharpe of
1.63. In contrast, look-ahead joins (e.g., using current-day closing prices before the
close) produce inflated paper Sharpe ratios up to 71.5. Abnormally high backtest performance
thus serves as an unambiguous signature of look-ahead leakage.

\subsection{Validation benchmarks}
\paragraph{Defect detection benchmark.}
We evaluate 13 core research guards against 12 planted quantitative defect families,
including look-ahead feature joins, unadjusted stock split prices, future earnings
disclosure leaks, and survivorship bias. Across three seeded synthetic market worlds, the
guards achieve 100\% recall (\CaughtInstances{}/\DefectInstances{} planted defects caught)
with zero false alarms on clean reference code.

\paragraph{Reporting discrepancy and compliance paradox.}
Table~\ref{tab:pilot} measures reporting discrepancy: the absolute difference between the
Sharpe ratio claimed by the agent in its written report and the ground-truth Sharpe ratio
computed by the oracle evaluator.

\begin{table}[t]
\centering\small
\begin{tabular}{lc}
\toprule
Model and condition & Mean absolute Sharpe gap \\
\midrule
Opus, no library & \OpusBase \\
Opus, with library & \OpusLibrary \\
Haiku, no library & \HaikuBase \\
Haiku, with library & \HaikuLibrary \\
\bottomrule
\end{tabular}
\caption{Reporting discrepancy between agent-claimed Sharpe ratios and ground-truth oracle evaluations.}
\label{tab:pilot}
\end{table}

Without tools, Opus exhibited a reporting gap of \OpusBase, which decreased to \OpusLibrary{}
when given execution guards (\OpusReduction\% reduction). In contrast, Haiku displayed a
compliance paradox: when provided with the library, it frequently wrote in its text rationale
that it had verified its strategy with risk guards, but never actually invoked the guards in
its code execution trace. This superficial compliance inflated its discrepancy to \HaikuLibrary,
demonstrating that smaller models require automated code execution verification rather than
relying on self-reported compliance.

\paragraph{Social sentiment verification.}
On 1.72 million financial posts across 2,521 social media accounts, raw uncalibrated sentiment
scores exhibited negative out-of-sample rank correlation with subsequent asset returns
($\text{Rank IC} = -0.0318$): blindly following social media sentiment lost money.
By evaluating each account's historical prediction accuracy and applying Bayesian credibility
weighting (including inverting signals from consistently inaccurate accounts), the out-of-sample
correlation flipped to $+0.0104$. This demonstrates that unstructured text signals require
empirical track-record verification before driving investment decisions.

\subsection{Four-condition audit feasibility}
We evaluated models under four audit conditions: unassisted baseline (C0), text guidance
only (C1), optional executable guards (C2), and mandatory automated audit (C3).

\paragraph{Operational findings.}
Table~\ref{tab:beacon_feasibility} presents results for Qwen2.5-Coder 14B under automated audit
constraints. The acceptance rate was 67\% under C0 and C1, 100\% under C2 (optional guards),
and 0\% under C3 (mandatory audit), where runs reached the turn limit while attempting to fix
all detected code defects. Mean token consumption under C3 was \BeaconFourteenCThreeMeanTokens{}
compared to \BeaconFourteenCTwoMeanTokens{} under C2, showing that automated verification
compels the agent to engage in substantially deeper code debugging before finalizing allocations.

\begin{table*}[t]
\centering\small
\begin{tabular}{llccc}
\toprule
Model & Condition & Accepted/Planned & Mean tokens & Primary status \\
\midrule
\multirow{4}{*}{14B}
  & C0: unassisted            & \BeaconFourteenCZeroAccepted/\BeaconFourteenCZeroPlanned
      & \BeaconFourteenCZeroMeanTokens   & 1$\times$ turn limit \\
  & C1: text guidance only    & \BeaconFourteenCOneAccepted/\BeaconFourteenCOnePlanned
      & \BeaconFourteenCOneMeanTokens    & 1$\times$ turn limit \\
  & C2: optional guards       & \BeaconFourteenCTwoAccepted/\BeaconFourteenCTwoPlanned
      & \BeaconFourteenCTwoMeanTokens    & all accepted \\
  & C3: mandatory audit       & \BeaconFourteenCThreeAccepted/\BeaconFourteenCThreePlanned
      & \BeaconFourteenCThreeMeanTokens  & 3$\times$ turn limit \\
\midrule
7B  & all (12 cells)          & 0/12 & --- & turn limit / syntax error \\
32B & all (12 cells)          & 0/12 & 0   & environment error \\
\bottomrule
\end{tabular}
\caption{Audit feasibility results across open-weight models under automated verification constraints.}
\label{tab:beacon_feasibility}
\end{table*}

A reference quantitative script passed all verification and timing checks on the market
benchmark, confirming the end-to-end viability of the automated verification pipeline.
